
@@M0@@
où @@M1@@ est la
dimension relative virtuelle de @@M2@@. Cet homomorphisme étant supposé défini, et tenant compte de la formule de Riemann--Roch sous la forme de Exp. VIII, à valeurs dans @@M3@@, on constate que la validité de la formule de Riemann--Roch @@M4@@-adique pour le morphisme @@M5@@ équivaut exactement à la commutativité du carré
@@M6@@
où les flèches horizontales se définissent à partir de (6.1) en passant aux gradués associés et en divisant, de façon à transformer classes de Chern en classes de Chern, et où les flèches verticales sont respectivement la flèche de Exp. VIII et l'homomorphisme de Gysin en cohomologie (supposé donné).

\textbf{6.2.} Notons que l'homomorphisme de Gysin est défini en tous cas si @@M7@@ est lisse sur un corps @@M8@@ (SGA 5 IV), en utilisant le théorème de pureté cohomologique relatif sur @@M9@@; la même construction s'applique lorsque @@M10@@ est excellent régulier et de caractéristique nulle, grâce au théorème de pureté cohomologique absolu de SGA 4 XIX, et elle s'appliquera automatiquement pour tout @@M11@@ excellent et régulier, sans restriction de caractéristique, une fois le théorème de pureté cohomologique absolu prouvé en toute généralité. Enfin, la même construction s'applique, grâce au théorème de pureté cohomologique relatif (SGA 4 XVI 3.7), lorsque @@M12@@ et @@M13@@ sont tous deux lisses sur un même schéma @@M14@@ et que @@M15@@ est un @@M16@@-morphisme. Il y a lieu de conjecturer que dans chacun de ces cas où l'homomorphisme de Gysin est défini, le diagramme (6.3) est bien commutatif. Le seul cas où cette assertion soit actuellement vérifiée est celui où @@M17@@ et @@M18@@ sont tous les deux lisses sur un corps @@M19@@, et résulte alors de la commutativité du diagramme analogue à (6.3) :
@@M20@@
